\scanpage{104}%
tions sur $F$ (décroissantes, \inserted{discrètes,} indexées par $\underline{Z}$) compatibles avec lesproduits tensoriels (en utilisant la notion év\deleted{\unsure{o|q}}\handwritten{i}dente de produit tensoriel\deleted{s} de deux espaces vectoriels filtrés). On notera qu'une telle donnée ne pourra plus s'exprimer par un homomorphisme d'un certain groupe algébrique dans $G$ (le groupe de Galois de $\underline{C}$ en $F$), car la catégorie des espaces vectoriels de dimension finie sur $k$, munis d'une filtration décroissante discrète indexée par $\underline{Z}$, n'est pas une catégorie abélienne (les bimorphismes ne sont pas des isomorphismes\handwritten{!}). Mais à une telle donnée est \unsure{associée|associé} un deuxième foncteur fibre\margin{\handwritten{$F'$}} $\inserted{\unsure{F|\mathbb{F}}'(V)=}\operatorname{Gr}F(V)$, à valeurs cette fois-ci dans les vectoriels gradués (NB $F$ joue le rôle de la cohomologie de \deleted{Hodge, \unsure{$F'$}} De Rham, $F'$ celui de la cohomologie de Hodge\deleted{\unsure{)}}, muni de la graduation \deleted{totale}\inserted{\handwritten{par $p$}}). La donnée de $F'$ correspond alors à un torseur \inserted{\handwritten{(à droite)}} $P$ sous $G$, le groupe de Galois correspondant $G'$ étant déduit de $G$ en tordant par $P$. La graduation de $F'$ correspond à un homomorphisme $i_{\handwritten{1}}:\underline{G}_m\handwritten{\longrightarrow}G'$. Considérons alors sur $F'$ la filtration décroissante associée à sa graduation, et soit $H'_{i_{\handwritten{1}}}=\underline{\operatorname{Aut}}\,\underline{\mathrm{filt}}^{\handwritten{!}}(F')\handwritten{\subset}\underline{\operatorname{Aut}}(F')=G'$ le sous-schéma en groupes de $G$ qui correspond aux \inserted{$\unsure{\otimes|\Omega}$-}automorphismes de $F'$ (ou plutôt des $F'_{k'}$, $k'$ une $k$-algèbre quelconque) qui respectent sa filtration\inserted{\handwritten{et induisant l'identité sur le gradué associé}}; il est canoniquement déterminé par $i_{\handwritten{1}}$. On peut aussi regarder le sous-schéma
\[
Q=\underline{\operatorname{Isom}}\,\underline{\mathrm{filt}}^{\handwritten{!}}(F,F')\handwritten{\subset}P=\underline{\operatorname{Isom}}(F,F')
\]
du schéma $P$ des $\otimes$-isomorphismes de $F$ avec $F'$, \deleted{formé de} qui correspond aux automorphismes respectant les filtrations de $F$ et $F'$\inserted{\handwritten{et induisant l'identité sur les gradués associés}}. C'est à priori un pseudo-torseur à \deleted{droite}\inserted{\handwritten{gauche}} sous $H'$, i.e. il est vide ou un torseur à droite sous $H'$. \uuline{Il faudrait prouver} que c'est bien un torseur (i.e. que sur une extension convenable $k'$ de $k$, on peut trouver un isomorphisme de $F_{k'}$ avec $F'_{k'}$ \uline{respectant les filtrations}). Donc on trouve une restriction du groupe d'opérateurs (à \deleted{droite}\inserted{\handwritten{gauche}}) $G'$ de $P$ au sous-groupe $H'$,
% archive sheet number: 103
